\section{Viabilidade econômica}

Não há base confiável para estimar preço, receita, custo de aquisição ou ponto de equilíbrio; por isso, esta versão não projeta faturamento. A avaliação começa com uma planilha de custos fixos mensais, taxas de pagamento, hospedagem, comunicação, atendimento e tempo da equipe, inclusive dos fundadores. O modelo da Atividade X lista assinaturas, planos para parceiros, intermediação e patrocínios como possibilidades [29]; aqui, elas são hipóteses independentes.

A receita mensal pode ser representada por
\[
R = (N_c \times P_c) + (N_a \times T_a) + S,
\]
em que $N_c$ é o número de clínicas pagantes, $P_c$ o preço do plano, $N_a$ os atendimentos intermediados, $T_a$ a receita líquida por atendimento e $S$ o patrocínio mensal. O ponto de equilíbrio em parceiros é
\[
N_{\text{eq}} = \frac{\text{custos fixos mensais}}{\text{margem de contribuição média por parceiro}}.
\]
O valor de um parceiro ao longo do tempo pode ser aproximado por $\text{LTV} = m \cdot \frac{r}{1 + d - r}$, em que $m$ é a margem por período, $r$ a taxa de retenção e $d$ a taxa de desconto [12]. A relação LTV/CAC só deve orientar investimento em aquisição quando houver dados reais. Os cálculos são métodos a aplicar após a coleta, não resultados desta atividade.

\begin{table}[H]
\centering\small
\caption{Hipóteses de receita e critérios de validação.}
\begin{tabularx}{\textwidth}{P{2.8cm}YY}
\toprule
\textbf{Hipótese} & \textbf{Como testar} & \textbf{Critério para avançar} \\
\midrule
Plano para parceiros & Entrevistas com clínicas sobre problema, valor e faixa de preço. & Duas clínicas aceitam novo ciclo e discutem preço. \\
Intermediação & Medir encaminhamentos, conclusão e custo em processo acompanhado. & Fluxo rastreável, custo conhecido e revisão jurídica. \\
Patrocínio & Proposta delimitada a empresas ou organizações locais. & Só contabilizar após manifestação formal. \\
Assinatura do tutor & Testar interesse sem bloquear funções essenciais. & Evidência de valor recorrente e intenção de pagamento. \\
\bottomrule
\end{tabularx}
\end{table}
